\section{Deeper thematic synthesis: where the methods agree and disagree}
\label{sec:deep-synthesis}

\subsection{The action interface is the hidden organizing variable}
The most consequential difference between several reviewed methods is not whether
they use a neural network. It is the meaning of the action they evaluate.
DWA evaluates planar velocities constrained by reachable and stoppable motion;
ORCA selects a collision-avoiding velocity under a simplified reciprocal model;
SARL's formulation makes the selected velocity instantaneous. Biped MPC instead
includes a balance-compatible stopping condition. At another level, a LiDAR-driven
locomotion policy emits joint commands, while a visual navigation model may emit
waypoints or trajectory-like targets for a separate controller
\citep{dwa,orca,sarl,biped2016,safehumanoid,flux}.

This difference changes the problem being solved. Let $\mathcal{A}_{\rm plan}$ be
the planner's nominal action set and $\mathcal{T}_{\rm exec}(u,x)$ the distribution
of physical trajectories produced by executing $u$ from state $x$. A geometric
planner effectively substitutes a deterministic nominal trajectory for this
distribution. A useful response-aware method must improve that substitution in
a way that affects decisions, not merely produce a better-looking predicted
path. In particular, the map from command to motion must include whatever
normalization, saturation, deadband, slew and policy history is active in the
deployed interface.

This framing also explains why post-hoc filtering is not innocuous. If a planner
checks a curved forward command and a supervisor replaces it with a stop, the
executed stopping trajectory may have different clearance from a moving person.
If a policy then produces lateral drift, neither previously checked path is the
real trajectory. A common command vocabulary is insufficient; the layers need
a common interpretation of the command's short-horizon physical effect. The
ORCA limitations and wheeled-legged failure analysis already motivate this
concern, so the novel object must be a specific shared model or coordination rule
\citep{orca,wheeledleg}.

\subsection{Reciprocity is an intervention, not a harmless simulator setting}
SARL's visible and invisible settings provide a particularly clear contrast.
The same general navigation problem produces very different ORCA performance
depending on whether the surrounding agents react to the robot. Its Table~I
reports success 0.43 for ORCA in the invisible setting, compared with 0.99 in
the visible setting in Table~II; these are distinct protocols rather than a
single algorithm failure rate \citep{sarl}. SocNavBench explicitly changes
collaboration coefficients for nonreactive pedestrians and describes the
kinematic choices used in baseline adaptation \citep{socnavbench}.

The experimental implication is stronger than ``use more realistic crowds''.
Reciprocal, independent and replayed actors answer different counterfactual
questions. Replay holds human trajectories fixed and can isolate robot-side
changes, but it removes the human response that an interactive planner is meant
to exploit. Reciprocal simulation permits that response but may make it too
predictable or overly accommodating. A data-driven actor can improve behavioral
diversity while introducing another learned model whose errors need evaluation.
One should therefore report results by crowd regime and avoid pooling them into
a universal social-navigation success rate.

For Geoff-Sim, the nonreactive regime is scientifically valuable precisely
because it tests unilateral responsibility. A contribution in that regime should
be named accordingly. Claims about cooperation, politeness or negotiation require
additional conditions in which the human can respond and the effects on the
human are actually measured \citep{survey2013,guidelines,proxemics2026}.

\subsection{A better representation can conceal a different experiment}
The transition from CADRL to SARL, RGL and DS-RNN is often summarized as a
sequence of more expressive networks. A closer reading shows simultaneous changes
in aggregation, dynamics modeling, memory, training and observation conditions.
RGL's comparison explicitly replaces an optimistic next-state assumption with a
linear model before attributing gains to relational prediction. DS-RNN introduces
a recurrent-attention ablation with matched PPO settings, as well as limited-view
and group conditions \citep{cadrl,sarl,rgl,dsrnn}. These decisions are part of
the contribution's evidential strength, not implementation trivia.

A proposed response-aware system must meet the same standard. If it receives
fresher tracks, a more capable policy, more CPU time and a different supervisor
at once, a higher completion rate cannot be assigned uniquely to calibration.
Conversely, running an externally trained baseline through an incompatible
action adapter can manufacture a failure. The comparison contract should list
inputs, predictor, action space, control rate, training data, runtime hardware,
post-processing and evaluator for every arm. A controlled ablation changes one
mechanism or uses a design that estimates interactions explicitly.

\subsection{Prediction quality has at least three distinct meanings}
First, a forecast can fit the recorded trajectory under an offline metric.
Second, it can represent uncertainty with useful probability or coverage.
Third, it can improve the decision made by the robot. Social LSTM, Social GAN
and Trajectron++ address different parts of this progression
\citep{sociallstm,socialgan,trajectron}. They cannot be collapsed into a claim
that any learned prediction is automatically safer than constant velocity.

Two details illustrate the danger. Social LSTM's Table~1 includes a privileged
destination input for IGP, and its aggregate final-displacement entry is not
better than every comparator. Social GAN evaluates the best of several samples
and explicitly discusses differences from the original Social LSTM training
setup. Trajectron++ adds distributional and dynamics-related metrics, including
KDE NLL, but probability calibration under deployment shift remains a separate
question. A method can generate a plausible future without assigning the right
probability to a collision-critical alternative.

For a practical planner, the relevant quantity is not just
$\|\hat p_h-p_h\|$ averaged over a test set. It is whether the predicted
occupied region $\mathcal{O}_h(\tau)$ contains the person at the times when
candidate robot trajectories approach them. Coverage should be examined by
horizon, occlusion, maneuver, range and observation age. A broad region can
increase coverage while making every route infeasible; the associated liveness
cost must be measured. This leads to a joint occupancy--execution problem rather
than a competition based on a single forecasting score.

\subsection{The 2016 biped paper changes the novelty boundary substantially}
The expanded primary reading of \citet{biped2016} is especially important.
Its formulation already includes center-of-mass/center-of-pressure dynamics,
foot-placement feasibility, a terminal capture/stopping condition and
uncertainty-expanded pedestrian occupancy. In its notation, a separation margin
can grow with position error and velocity error over time. That is direct prior
art against a claim that combining biped dynamics, uncertainty and stopping is
new.

The paper also distinguishes passive safety from actually preventing every
collision. Its numerical people do not avoid the robot, and the study varies
field of view, horizon, crowd speed and uncertainty. A longer preview is not
uniformly helpful when the sensor cannot observe the relevant closing region.
Most strikingly, one relaxation reduces collisions partly by allowing infeasible
robot motion and falls; another preserves dynamic feasibility and obtains a
better total-failure tradeoff. This is a concrete demonstration that optimizing
one safety proxy can transfer risk to another outcome.

The defensible extension is therefore narrower: identify and validate the
short-horizon response of a deployed learned locomotion policy, connect it to
acquisition-time uncertainty under real onboard perception, and test its value
relative to a generic uncertainty-aware/biped-feasibility model. The distinction
between a reduced-order analytic walking model and an empirically calibrated
policy interface could matter, but it is an experimental hypothesis, not a
novelty result established by changing robot names.

\subsection{Read recent end-to-end safety claims at the level of constraints}
The safe-humanoid paper presents a formal definition of safe behavior and a
barrier-inspired training mechanism. These are not the same object. Its text
states that the barrier violation is used as a cost rather than as an online
action constraint \citep{safehumanoid}. Approximate optimization of an expected
discounted cost does not directly provide the pointwise invariance property of
an online feasible CBF controller \citep{cbf,safecontrol}. This distinction is
compatible with useful empirical performance; it simply prevents an unjustified
guarantee from being carried across interfaces.

The tables likewise support a qualified interpretation. The proposed method has
better dynamic-agent and narrow-passage results in the reported comparison, but
not the best suspended-obstacle result. Nonzero unsafe-region duration remains,
and the evaluation and training distance thresholds differ. The ablation prose
also makes it difficult to isolate every contribution of the barrier-based cost
from the comfort-related rewards without a fully factorial comparison. These
are specific reasons to request matched variants and definitions, rather than
to dismiss constrained RL as a family.

Omni-Perception adds another lesson. Its PD-RiskNet ablation reports both
76.7\% success and 56.7\% collision for one method; those are evidently not
complementary exclusive outcome categories in the printed table
\citep{omni}. Without the event semantics, success cannot be reinterpreted as
collision-free completion. The paper's qualitative sensor comparison, rendering
throughput and real-world scenario table are also separate evidence products.
They should remain separate in an analysis of sensing fidelity, computational
efficiency and locomotion reliability.

\subsection{General-purpose navigation broadens the task but not automatically the guarantee}
FLUX, OmniNavBench and DPed-VLN extend the field toward dynamic scenes,
cross-embodiment deployment and composed or language-conditioned tasks
\citep{flux,omninav,dped}. They demonstrate that the ambition of the project
cannot be positioned as the first movement beyond isolated point avoidance.
Their evaluation contracts, however, differ substantially.

FLUX defines a collision-related statistic as intimate-space intrusion and uses
task termination based on progress, timeout and goal proximity. Its inference
timing and remote deployment setup are not comparable with a local-planner
callback p95. OmniNavBench distinguishes subtask predicates, complete-episode
success and personal-space intrusion over time; strong partial progress can
coexist with poor whole-mission success. DPed-VLN shows that a method can improve
goal-reaching while worsening a human-collision metric, and that adding moving
pedestrians can change risk much more than it changes success.

These comparisons motivate an explicit metric ontology: geometric overlap,
native contact, intimate-space intrusion, personal-space duration, failure to
stop, fall and incomplete mission are separate events. A common name such as
``collision rate'' does not make the events equivalent. A future benchmark
contribution should publish the event predicates and time basis so that readers
can reconstruct exactly what was measured.

\subsection{Perception and execution uncertainty interact rather than simply add}
DROW and DR-SPAAM show that surface measurements need normalization and learned
aggregation to estimate an object's center; JRDB shows that misses, localization
error and identity continuity can have different rankings
\citep{drow,drspaam,jrdb}. FAST-LIO2 and ORB-SLAM3 show that robot-state
estimation itself has a sensor, timing and calibration contract
\citep{fastlio,orbslam}. A pedestrian position transformed through an erroneous
robot pose is not an independent error source from the robot's estimated route.

An uncertainty-aware navigation study should therefore separate measurement
noise, systematic center bias, unobserved motion, state-estimation error and
command-response residuals. A single isotropic margin can conceal these
mechanisms. Nonetheless, adding a separate learned distribution for every error
source may make inference slow enough to worsen the total problem. The objective
is not maximal model complexity, but a representation that changes the right
decisions at the available observation age and compute budget.

This is why a source-audited factorial experiment is valuable. Replace one
estimated component at a time with its evaluation counterpart, then test
interactions between factors. If a better response model helps only when person
observations are time-aligned, that interaction is a substantive result. If it
does not help even with oracle perception, the main limitation may be the
controller or decision rule rather than the detector.

\subsection{Embodied adaptation makes the ``same controller'' a qualified statement}
RMA, perceptive locomotion, ExBody and ASAP demonstrate different ways in which
experience, sensor belief, posture and dynamics alignment affect physical
behavior \citep{rma,rough,expressive,asap}. A network file and a maximum speed
do not fully identify the executed dynamics. Observation history, command
normalization, PD gains, arm authority, payload and reset state matter as well.
The wheeled-legged study's tracking and memory failures reinforce that this is
already an active systems-level research problem \citep{wheeledleg}.

For social locomomanipulation, a reaching arm changes both dynamics and
collision geometry. A base-only envelope may be a useful approximation, but it
requires validation across arm poses and load transitions. First measure where
the approximation changes collision decisions. Only then is the extra cost of
articulated swept-volume checking justified. A failure of the simpler model can
be a meaningful finding; a full-body model that adds latency without improving
outcomes is not automatically a better contribution.

\subsection{The social-validity gap is contextual, not the absence of human research}
SCAND includes human evaluation, Habitat~3.0 supports human-in-the-loop
interaction, and the guidelines/proxemics work explicitly studies how social
claims should be evaluated \citep{scand,habitat3,guidelines,proxemics2026}.
It would therefore be incorrect to motivate the project by saying that existing
social navigation has never considered people themselves.

The narrower question is whether behavior that improves the measured
collision--completion tradeoff on a humanoid is also experienced as more
predictable, comfortable and appropriate in a declared task. Morphology, sensor
height, arm motion, carried objects and stopping style can all affect that
relationship. An informative study would control speed and clearance while
varying the response-aware decision mechanism, use counterbalanced conditions,
and report both objective outcomes and participant-level judgments. The results
would then support a bounded statement about that task and population, not a
universal preferred personal-space shape.

\subsection{What this synthesis supports}
Across the themes, the most defensible opening is an \emph{auditable coupling}
between observation age, uncertainty, executable humanoid response and the
safety--liveness decision. Existing work establishes the individual ingredients
and several combinations. The project must therefore demonstrate either a
specific new mechanism with an advantage under matched conditions, or a new
controlled explanation of failure mechanisms that changes how systems should
be designed. The next section's gap priorities and the final novelty roadmap
state the evidence needed to move from that opening to a contribution claim.
